\section{Finite exact certificate completeness}
\label{sec:certificates-independent}

Fix one principal-set branch of the semialgebraic characterization. Replace
every strict inequality by an equation:
\[
 g>0\iff\exists z\ (gz^2-1=0),\qquad
 g<0\iff\exists z\ ((-g)z^2-1=0).
 \tag{EQ}
\]
These equivalences are exact over the reals: a solution forces the indicated
sign, and a strict sign supplies $z=1/\sqrt{|g|}$.

\begin{theorem}[Two-sided finite certificates]
\label{thm:certificates-independent}
Every YES instance of the exact all-mobile decision problem has a finite
real-algebraic sample certificate. Every NO instance has a finite family of
Real Nullstellensatz identities, one for each nonempty principal branch (or,
equivalently, one identity for the Boolean-selector encoding). Both kinds are
independently checkable by exact arithmetic. No polynomial bound on their size
is asserted.
\end{theorem}

\begin{proof}
A nonempty semialgebraic set defined over $\mathbb Q$ contains a point with
real-algebraic coordinates. The finitely many coordinates generate a finite
real extension of $\mathbb Q$, which is simple: represent it as
$\mathbb Q(\alpha)$ using the irreducible minimal polynomial of $\alpha$, a
rational isolating interval selecting the real embedding, and polynomial
representatives for every coordinate. Exact reduction modulo the minimal
polynomial checks all equations; Sturm or subresultant sign determination at
the isolated root checks every required sign. This gives a finite YES
certificate tied to the reconstructed network formula rather than to an
unresolved symbolic parameter system.

For an infeasible branch, let $F_1,\ldots,F_N$ be the polynomial equations
after (EQ). The Real Nullstellensatz states that emptiness of their real zero
set is equivalent to the existence of polynomials $q_k$ and $s_j$ with
\[
 -1=\sum_j s_j^2+\sum_{k=1}^Nq_kF_k.
 \tag{RN}
\]
A verifier expands both sides and compares every rational coefficient. At a
hypothetical common real zero, (RN) would give
$-1=\sum_j s_j^2\geq0$, proving soundness. There are finitely many principal
branches, so a NO instance has a finite collection of such identities.
\end{proof}

Existence, verification, and generation must be distinguished. Identities and
algebraic samples can in principle be generated by exhaustive enumeration of
degrees, rational coefficients, minimal polynomials, isolating intervals, and
exact identity checks; termination follows from the theorem. The result does
not provide a polynomial certificate-size bound or a practical general
certificate generator. A bundled classifier may therefore return
\texttt{UNRESOLVED} on a raw network without contradicting the decision or
certificate-existence theorems.
